\section{运行状态划分与 Markov 依赖}

\subsection{为何需要状态框架}

如果只使用全样本频率与严重度分布，那么模型默认所有月份来自同一种环境。这与实际运营并不一致：机场存在正常月份，也存在明显高压、恢复困难、取消增加的扰动月份。因此，项目先在机场月份层面定义状态，再讨论状态之间的转移。

\subsection{候选状态规则}

项目比较了两套状态规则：
\begin{enumerate}[leftmargin=2.2em]
  \item \textbf{impact\_q75}：如果月份总冲击高于 \texttt{disruption\_impact\_minutes} 的 75\% 分位数，则记为 disrupted；
  \item \textbf{delay\_cancel\_impact\_composite}：把延误率、取消率和 delay-equivalent minutes 标准化后求和，再取 75\% 分位以上为 disrupted。
\end{enumerate}

第二种规则的数学形式为：
\begin{align}
z_t^{(d)} &= \frac{\text{delay\_rate}_t - \mu_d}{\sigma_d}, \\
z_t^{(c)} &= \frac{\text{cancellation\_rate}_t - \mu_c}{\sigma_c}, \\
z_t^{(L)} &= \frac{L_t - \mu_L}{\sigma_L}, \\
\text{Composite}_t &= z_t^{(d)} + z_t^{(c)} + z_t^{(L)}.
\end{align}

然后令
\begin{equation}
\text{state}_t =
\begin{cases}
\text{disrupted}, & \text{if } \text{Composite}_t \ge Q_{0.75}(\text{Composite}), \\
\text{normal}, & \text{otherwise}.
\end{cases}
\end{equation}

\subsection{规则选择方法}

项目没有主观指定状态规则，而是构造了一个 selection score：
\begin{equation}
\text{score}
= \left| \text{disrupted\_share} - 0.25 \right|
+ \operatorname{std}\left(\text{yearly disrupted share}\right).
\end{equation}

第一项控制 disrupted 月份占比不要偏离 25\%，第二项控制不同年份之间 disrupted 占比的稳定性。分数越低，说明状态定义既保留了 ``四分之一高风险月份'' 的直觉，又避免年份之间波动过大。

结果如下：
\begin{table}[H]
\centering
\caption{状态规则比较}
\begin{tabular}{lrrrrc}
\toprule
规则 & disrupted 占比 & 年度占比标准差 & disrupted 月数 & score & 是否选中 \\
\midrule
delay\_cancel\_impact\_composite & 0.25 & 0.1521 & 18 & 0.1521 & Yes \\
impact\_q75 & 0.25 & 0.1863 & 18 & 0.1863 & No \\
\bottomrule
\end{tabular}
\end{table}

因此，项目最终选择 \texttt{delay\_cancel\_impact\_composite} 作为状态识别规则。

\begin{figure}[H]
  \centering
  \includegraphics[width=0.82\textwidth]{chart_10_state_rule_comparison_en.png}
  \caption{状态规则比较图}
  \label{fig:state-rule}
\end{figure}

\subsection{状态划分后的经验特征}

在最终选中的规则下：
\begin{itemize}[leftmargin=2.2em]
  \item 72 个机场月份中，54 个为 normal，18 个为 disrupted；
  \item normal 状态下，平均月度冲击为 149{,}550.65 分钟；
  \item disrupted 状态下，平均月度冲击为 376{,}723.50 分钟；
  \item normal 状态平均取消率为 0.0147，而 disrupted 状态提高到 0.0811。
\end{itemize}

\begin{figure}[H]
  \centering
  \includegraphics[width=0.78\textwidth]{chart_4_disruption_state_impact.png}
  \caption{不同状态下的航空公司月份冲击分布}
  \label{fig:state-impact}
\end{figure}

\begin{figure}[H]
  \centering
  \includegraphics[width=0.95\textwidth]{chart_5_monthly_disrupted_share.png}
  \caption{机场月份状态时间线}
  \label{fig:state-timeline}
\end{figure}

样本中的峰值扰动月份是 2023 年 7 月，其 delay-equivalent minutes 达到 583{,}999 分钟。随后较高风险月份还包括 2025 年 7 月、2022 年 6 月、2024 年 8 月等。图 \ref{fig:state-timeline} 说明高风险月份并非随机散点，而是呈现出季节性与阶段性聚集。

\subsection{Markov 转移矩阵的估计}

在状态被定义后，项目按月份时间顺序构造一阶 Markov 链。设状态空间为
\begin{equation}
\mathcal{S} = \{\text{normal}, \text{disrupted}\}.
\end{equation}

对任意 $i,j \in \mathcal{S}$，转移概率估计为
\begin{equation}
\hat{p}_{ij} = \frac{N_{ij}}{\sum_{k \in \mathcal{S}}N_{ik}},
\end{equation}
其中 $N_{ij}$ 表示从状态 $i$ 转移到状态 $j$ 的观察次数。

当前样本得到的转移矩阵为
\begin{equation}
\hat{P}
=
\begin{bmatrix}
0.814815 & 0.185185 \\
0.529412 & 0.470588
\end{bmatrix},
\label{eq:transition-matrix}
\end{equation}
其中行对应当前状态，列对应下一期状态，即
\begin{equation}
\hat{P} =
\begin{bmatrix}
\mathbb{P}(N \to N) & \mathbb{P}(N \to D) \\
\mathbb{P}(D \to N) & \mathbb{P}(D \to D)
\end{bmatrix}.
\end{equation}

对应的转移计数是：
\begin{table}[H]
\centering
\caption{Markov 转移计数与概率}
\begin{tabular}{llrr}
\toprule
from & to & 计数 & 概率 \\
\midrule
normal & normal & 44 & 0.8148 \\
normal & disrupted & 10 & 0.1852 \\
disrupted & normal & 9 & 0.5294 \\
disrupted & disrupted & 8 & 0.4706 \\
\bottomrule
\end{tabular}
\end{table}

\begin{figure}[H]
  \centering
  \includegraphics[width=0.55\textwidth]{chart_7_markov_transition_matrix.png}
  \caption{两状态 Markov 转移矩阵}
  \label{fig:markov}
\end{figure}

\subsection{平稳分布的计算}

若记平稳分布为 $\bm{\pi} = (\pi_N,\pi_D)$，则有
\begin{equation}
\bm{\pi}\hat{P} = \bm{\pi}, \qquad \pi_N + \pi_D = 1.
\end{equation}

对二状态链，可直接解得
\begin{equation}
\pi_D = \frac{\hat{p}_{ND}}{\hat{p}_{ND} + \hat{p}_{DN}}
      = \frac{0.185185}{0.185185 + 0.529412}
      \approx 0.2591,
\end{equation}
以及
\begin{equation}
\pi_N = 1 - \pi_D \approx 0.7409.
\end{equation}

这说明在长期意义上，系统大约有 74.1\% 的时间停留在 normal 状态，有 25.9\% 的时间停留在 disrupted 状态。这个结果与前面状态规则中 ``约四分之一月份为 disrupted'' 的设计逻辑是一致的。

\subsection{本阶段的解释}

Markov 结果可以读成一种非常直接的运营叙事：
\begin{itemize}[leftmargin=2.2em]
  \item 正常状态具有明显惯性，下一月继续保持 normal 的概率约为 81.5\%；
  \item 一旦进入 disrupted 状态，仍有 47.1\% 的概率在下一月继续维持 disrupted；
  \item 因此，扰动不仅是高损失月份本身，更体现为 \textbf{状态持续性与恢复难度}。
\end{itemize}
